\subsection{Extended EA Addressing Modes}\label{section:extended-ea-addressing}
EXT1 and EXT2 extend EA when a compact form cannot express the operand: explicit segment selection, indexed addressing,
zero-base addressing, SP/PC indexed addressing, or auto-update addressing. EXT1 selects exactly one descriptor byte;
EXT2 selects exactly two descriptor bytes.

Each extended-descriptor form selects its segment from the encoded SREG (:term:base.terms.field:) SEG(s), the fixed SS segment for SP, the fixed CS
segment for PC, or the operation\textquotesingle{}s default data segment.

EA evaluation begins by copying the architectural registers into (:term:base.terms.temporary_operand_evaluation_state:). Operands are
then processed in instruction order. Within one EA, the base term is evaluated before the index term and then the
displacement. Auto-update changes only the (:term:base.terms.temporary_operand_evaluation_state:) while operands are being evaluated. Memory reads and
writes are collected as instruction side effects, and neither those effects nor the register updates become
architecturally visible until the instruction commits. A fault before commit therefore leaves both register and
memory state unchanged apart from the exception-entry transition.

Postincrement forms use the current temporary value and update it afterward. Predecrement forms update the temporary
value first and then use the result. A base register changes by the EA interpretation width in bytes; an index
register changes by one element before scaling. Each REP iteration applies and commits these rules independently.
Explicit register references observe preceding updates in the temporary evaluation state. These EA evaluation-order
rules also apply when a field-selected register and an explicit register reference name the same register.

Size-less address-role instructions INVPAGE, FLSHDCACHE, INVDCACHE, INVICACHE, WRBKDCACHE, SYNCCACHE, PREFETCH,
PREFETCHNT, and PTQUERY have Q interpretation width. Size-less control-target instructions DJcc, IJcc,
LCALL, and LJMP also have Q interpretation width. Their extended-descriptor index scale is eight and their base predecrement or
postincrement amount is eight. Cache range loops advance their explicit address by the CPUID maintenance granule,
independently of this per-EA update amount.

When one register supplies both terms in \texttt{[R1 + R1++ * scale]}, the base and index terms use the same
current temporary R1 before the index update. In \texttt{[R1 + {-}{-}R1 * scale]}, the base uses the current value;
the index term then decrements temporary R1 and uses the updated value.

\paragraph{Scalar EXT1 Modes}\mbox{}\par
(:ea-diagram:base.ea.modes.EXT1.explicit_segment_base:)
\clearpage
(:ea-diagram:base.ea.modes.EXT1.explicit_segment_zero_base:)
\clearpage
(:ea-diagram:base.ea.modes.EXT1.default_segment_base:)

\clearpage
\paragraph{Scalar EXT2 Modes}\mbox{}\par
(:ea-diagram:base.ea.modes.EXT2.explicit_segment_index:)
\clearpage
(:ea-diagram:base.ea.modes.EXT2.explicit_segment_base:)
\clearpage
(:ea-diagram:base.ea.modes.EXT2.explicit_segment_zero_base_index:)
\clearpage
(:ea-diagram:base.ea.modes.EXT2.stack_pointer_index:)
\clearpage
(:ea-diagram:base.ea.modes.EXT2.program_counter_index:)
